\clearpage
\section*{Public Blog: How Much Comfort Does One Kilowatt-Hour Buy?}
\addcontentsline{toc}{section}{Public blog}

An outdoor fan can help people feel cooler while the air thermometer barely changes. An evaporative cooler lowers air temperature while adding moisture. A fair comparison follows the cooling to the visitor: what temperature, humidity, air movement and radiant heat reach that person, and for how long?

Our model combines those influences using Standard Effective Temperature, or SET. One service unit means reducing one person's SET excess by one degree for one hour. Dividing that benefit by electricity gives cooling efficiency. It measures relief accumulated over visitor-hours.

We applied the same calculation to two very different spaces: a large occupied corridor under the Fremont Street Experience canopy in Las Vegas, and a shaded ticket queue at Honolulu Zoo \cite{fremont-canopy,honolulu-zoo-entrance}. The modeled areas are 8382 and 60 square meters. Las Vegas offers dry heat; Honolulu offers a warm, humid coastal setting. The smaller queue also packs more visitor-time into each square meter, which helps each unit serve more people.

At Honolulu, three fans cover the queue and deliver about \textbf{143 service units using 3 kWh}: roughly \textbf{48 units per kWh}. In Fremont, placing one fan in the busiest demand zone yields about \textbf{17 units per kWh}, but directly serves only \textbf{0.33\% of visitor-time} across the corridor. A good efficiency ratio can coexist with very limited coverage. More equipment in the same busy zone can retain that ratio while serving more people.

Choosing a practical installation adds another question: how much service must it provide, and at what cost? Our illustrative policy requires at least 80\% coverage and $1.5^\circ$C average relief, alongside space and electricity limits. It selects 266 fans plus 16 evaporative units in Fremont, and one fan plus one evaporative unit in Honolulu. In the queue, three fans already deliver $1.4908^\circ$C relief. Lowering the target from $1.50$ to $1.49^\circ$C changes the least-cost answer to fans alone.

Scale brings a similar distinction. If Honolulu's area and crowd both grow eightfold, keeping three fans preserves their efficiency but lowers coverage from 100\% to 12.5\%. Twenty-four fans restore coverage. The useful questions are therefore connected: how much benefit does each kWh buy, who receives it, and what service can the operator afford? During severe heat, also follow local advisories and adjust outdoor plans \cite{nws-heat}.
